\subsection{Die Cramersche Regel}
\label{subsec:linear_systems_cramers_rule}

Die Formeln für $x_1$ und $x_2$ besitzen ein auffällige Struktur:
\begin{itemize}
  \item Im Nenner steht jeweils die Determinante der Koeffizientenmatrix.
  \item Im Zähler steht die Determinante einer Matrix, in der die zur
  gesuchten Unbekannten gehörende Spalte durch die Absolutglieder ersetzt wurde.
\end{itemize}
Für eine kompakte Schreibweise führen wir dafür die Cramerschen Ersatzmatrizen ein.

\begin{DefinitionBox}{Cramersche Ersatzmatrix}{linear_systems_cramer_replacement_matrix}
  Gegeben sei ein $n\times n$-LGS mit Koeffizientenmatrix $A$ und dem Vektor
  der Absolutglieder $b$. Ersetzt man in $A$ die $i$-te Spalte durch $b$, so
  entsteht die $i$-te \emph{Cramersche Ersatzmatrix}
  \EnglishTerm{Cramer replacement matrix} $A_i$.
  \par\medskip
  Beispielhaft für den $2\times2$-Fall:
  \begin{center}
    \CachedTikzFigure{content/resources/chapters/linear_systems/tikz/cramers_column_replacement.tex}
  \end{center}
\end{DefinitionBox}

Mit den Erstzmatrizen $A_i$ lässt sich die Lösung des allgemeinen $2 \times 2$ LGS
nochmals kompakter und übersichtlicher formulieren.
%
\begin{samepage}
Für ein $2\times2$-LGS mit $\det(A)\neq0$ lauten die Lösungen damit
\[
  x_1=\frac{\det(A_1)}{\det(A)},
  \qquad
  x_2=\frac{\det(A_2)}{\det(A)}.
\]

Die Cramersche Regel verallgemeinert dieses Prinzip auf beliebig viele Unbekannte.
\begin{TheoremBox}{Cramersche Regel}{linear_systems_cramers_rule_two_by_two}
  Für ein $n\times n$-LGS mit Koeffizientenmatrix $A$ und $\det(A)\neq0$
  gilt
  \[
    x_i=\frac{\det(A_i)}{\det(A)},
    \qquad i=1,\dots,n.
  \]
  Diese Lösungsformel heißt \emph{Cramersche Regel}
  \EnglishTerm{Cramer's rule}.
\end{TheoremBox}
\end{samepage}

\begin{ExampleBox}{Gauß-Verfahren versus Cramersche Regel}{linear_systems_cramers_rule_example}
  Wir lösen das folgende LGS mit 2 verschiedenen Verfahren.
  \[
    \begin{aligned}
      x_1+x_2&=2,\\
      x_1-x_2&=-4.
    \end{aligned}
    \qquad\Longleftrightarrow\qquad
    [A\mid b]=
    \left[
      \begin{array}{@{}rr|r@{}}
        1& 1& 2\\
        1&-1&-4
      \end{array}
    \right]
  \]
  {\setlength{\tabcolsep}{0.1cm}%
  \setlength{\arrayrulewidth}{0.7pt}%
  \noindent\begin{tabularx}{\linewidth}{@{}X|X@{}}
  \centering\textbf{Gauß-Verfahren}
  &
  \centering\arraybackslash\textbf{Cramersche Regel}\\[0.2em]
  \hline
  \begin{minipage}[t]{\linewidth}
    \vspace{0.45em}
    \begin{ClozeArea}[16]{22}
      Vorwärtselimination:
      \[
        \begin{array}{r@{\;}c@{\qquad}l}
          &\left[\begin{array}{@{}rr|r@{}}
            1& 1& 2\\
            1&-1&-4
          \end{array}\right]
          &\raisebox{-0.65em}{\(\TransformOp{Z_2-Z_1}\)}\\[1.0em]
          \Longleftrightarrow
          &\left[\begin{array}{@{}rr|r@{}}
            1& 1& 2\\
            0&-2&-6
          \end{array}\right]
          &\text{Zeilenstufenform}
        \end{array}
      \]
      \smallskip
      Rückwärtseinsetzen:
      \begin{itemize}[leftmargin=1.4em,itemsep=0.3em,topsep=0.2em]
        \item In Zeile 2 einsetzen:
          \[
            -2x_2=-6
            \quad\Longleftrightarrow\quad x_2=3.
          \]
        \item In Zeile 1 einsetzen:
          \[
            \begin{aligned}
              x_1+x_2=2
              &\quad\Longleftrightarrow\quad x_1+3=2,\\
              &\quad\Longleftrightarrow\quad x_1=-1.
            \end{aligned}
          \]
      \end{itemize}
      Lösungsmenge:
      \[
        \mathcal L=\{(-1,3)\}.
      \]
    \end{ClozeArea}
  \end{minipage}
  &
  \begin{minipage}[t]{\linewidth}
    \vspace{0.45em}
    \begin{ClozeArea}[16]{22}
      \[
        A=\begin{pmatrix}1&1\\1&-1\end{pmatrix}.
      \]
      Cramersche Ersatzmatrizen:
      \[
        A_1=\begin{pmatrix}2&1\\-4&-1\end{pmatrix},
        \qquad
        A_2=\begin{pmatrix}1&2\\1&-4\end{pmatrix}.
      \]
      Determinanten:
      \[
        \begin{aligned}
          \det(A)
            &=\det\!\begin{pmatrix}1&1\\1&-1\end{pmatrix}
              =1\cdot(-1)-1\cdot1=-2,\\[0.8em]
          \det(A_1)
            &=\det\!\begin{pmatrix}2&1\\-4&-1\end{pmatrix}
              =2\cdot(-1)-1\cdot(-4)=2,\\[0.8em]
          \det(A_2)
            &=\det\!\begin{pmatrix}1&2\\1&-4\end{pmatrix}
              =1\cdot(-4)-2\cdot1=-6.
        \end{aligned}
      \]
      Cramersche Regel:
      \[
        x_1=\frac{\det(A_1)}{\det(A)}
        =\frac{2}{-2}=-1,
      \]
      \[
        x_2=\frac{\det(A_2)}{\det(A)}
        =\frac{-6}{-2}=3.
      \]
      Lösungsmenge:
      \[
        \mathcal L=\{(-1,3)\}.
      \]
    \end{ClozeArea}
  \end{minipage}
  \end{tabularx}}
  \par\smallskip
  Mit beiden Verfahren wird dieselbe Lösungsmenge ermittelt.
\end{ExampleBox}

Um die Cramersche Regel auch für allgemeine $n \times n$ lineare Gleichungssysteme
anwenden zu können, müssen wir zunächst lernen, wie eine $n \times n$-Determinante
berechnet wird.
